\documentclass[12pt]{article}

\usepackage{geometry}
\usepackage[T1]{fontenc}
\usepackage[utf8]{inputenc}
\usepackage[frenchb]{babel}

\usepackage{graphicx}
\usepackage{lmodern}

\geometry{margin=2cm}


\begin{document}

%%%%%%%%%%%%%%%%%%%%%%%%%%%%%%%%%%%%%%%%%%%
\thispagestyle{empty}
\noindent\includegraphics[width=0.25\textwidth]{enseirb-matmeca}

\vspace{\stretch{1}}

\begin{center}
    {\Huge\textbf{Projet de compilation}

\vspace{\stretch{1}}

    \textbf{Rubic}}
\end{center}

\vspace{\stretch{2}}

\begin{tabular}{r@{:~}l}
	\textbf{Auteurs} & \textit{David Bitonneau, Benoît Ruelle}\\
	\textbf{Encadrant} & \textit{M. Denis Barthou}\\
\end{tabular}

\vspace{\stretch{1}}

\begin{center}Deuxième année, filière informatique

    Date : \today
\end{center}

\newpage
%%%%%%%%%%%%%%%%%%%%%%%%%%%%%%%%%%%%%%%%%%%

\section{Introduction}

Ce projet a pour but de créer un compilateur pour le langage Rubic, un dérivé
simplifié du langage Ruby. Le code généré n'est pas de l'assembleur mais un
langage intermédiaire le LLVM, assez proche du langage machine mais ayant ses
spécificités et étant notamment indépendant de tout système. Le code généré
est ainsi portable et peut ensuite être traité par le compilateur LLVM qui
pourra faire des optimisations lors de la compilation vers un langage machine.

L'accomplissement du projet a nécessité plusieurs étapes : la détermination de
la grammaire du langage Rubic, le typage des données et la génération du code
LLVM.
Afin de faciliter la réalisation de ce compilateur, les outils d'analyse
lexicale et syntaxique Flex (version libre de Lex) et Yacc ont été utilisés.

\section{Langage Rubic}

\subsection{Les variables}

En Ruby les noms de variables sont des étiquettes sans type et ce sont les
objets qu'elles contiennent qui sont fortement et dynamiquement typé. Il n'y a
donc pas de vérification de typage statique et c'est à l'exécution qu'il est
vérifié si l'appel de méthode fonctionne pour l'objet traité. Une même
variable peut donc contenir des objets de types différents au cours de sa vie.

En Rubic, le typage est également dynamique et une même variable peut
contenir différents types d'objets au cours de sa vie, cependant les
vérifications sont effectuées à la compilation et non à l'exécution.

\paragraph{}
En Rubic il existe différents types de variables :
\begin{itemize}
    \item les nombres entiers ;
    \item les nombres flottants ;
    \item les booléens ;
    \item les chaînes de caractères ;
    \item les objets\footnote{En Ruby tout est un objet, même les classes et
        les entiers.}.
\end{itemize}

\paragraph{}
Ces variables peuvent avoir différentes portées :
\begin{itemize}
    \item Les variables qui commencent par un '\$' sont des variables
        globales, elles sont accessibles partout une fois qu'elles ont été
        définies.
    \item Les variables qui commencent par un '@' sont des variables
        d'instances, leur portée est limitée à l'objet et elles ne peuvent
        être accédées à l'extérieur que via des méthodes.
    \item Les autres variables sont locales aux environnements auxquelles
        elles appartiennent. Par exemple elles peuvent être locale à la
        fonction dans laquelle elles ont été définies ou si elles ont été
        définies à l'extérieur de toute fonction ou classe, elles ne peuvent
        être accédées à l'intérieur de ces dernière.
\end{itemize}

\paragraph{}
Enfin, par convention, les variables dont le nom commence par une majuscule
sont des constantes. En Ruby, leur modification entraine seulement l'affichage
d'un avertissement et est effective. \textbf{[En Rubic, nous avons décidé de
déclencher une erreur de compilation lorsqu'une constante est modifiée, ce qui
signifie que le processus de compilation ne s'achèvera pas.{\Large/}En
Rubic nous avons pris la même décision, seul une message d'avertissement est
affiché lorsque des constantes sont modifiées et la compilation s'achèvera
normalement.}]\emph{(rayer la mention inutile)}

\subsection{Les fonctions}

\subsection{Les classes}

\section{Analyse sémantique}

\subsection{Tables de hashage}

Notre compilateur propose trois espaces de noms : un pour les
variables, un autre pour les fonctions et un dernier pour les classes. Cela
permet à une variable de coexister avec une classe ou une fonction du même
nom. Chaque espace de nom est implémenté à l'aide d'une table de hashage dont
le fonctionnement est le suivant.

Deux éléments sont demandés : un identificateur et un poiteur vers une donnée.
Une clé est obtenue en sommant tous les octets de l'identificateur modulo la
taille de la table de hashage (constante fixée à la compilation). Pour prendre
en charge les éventuelles collisions, à chaque case de la table correspond une
liste chaînée (une LIFO pour réutiliser ce qui a déjà été implémenté) où sont
insérés les données ayant la même clé (eg. "toto" et "otto" produisent la même
clé) et les identifiants avant le hashage sont alors conservés dans la liste
pour distinguer deux éléments lors de la recherche. La recherche est alors
simple : à partir de l'identifiant recherché, reproduire la clé et parcourir
la liste dans la case de la table de hashage correspondante.


\subsection{Visibilité des symboles}

La visibilité des symboles a été prise en charge à l'aide d'une pile et d'une
structure appelée "block" comprenant trois tables de hashage (une par espace
de nom). Un nouveau bloc est empilé à chaque fois qu'un lexème définissant une
nouvelle visibilité est lu. C'est le cas par exemple de la règle IF/THEN/ELSE
:

\begin{verbatim}
IF
    { nouveau bloc pour IF afin de stocker les éléments de expr
      qui doit rester visible dans le THEN et dans le ELSE }
expr
THEN
    { nouveau bloc pour stocker les éléments du bloc IF }
stmts
terms
ELSE
    { dépiler le bloc THEN et créer un nouveau bloc pour le ELSE.
      ici le bloc du IF contenant les éléments de expr est tojours
      visible. }
stmts
terms
END
    { dépiler le bloc ELSE et celui du IF }
\end{verbatim}


La recherche d'un symbole se fait en parcourant la pile pour
tester sa présence dans la table de hashage du type correspondant dans chaque
bloc.

\section{Typage}

\section{Génération de code}

\section{Conclusion}

Les exigences du projet ont été remplies et des améliorations ont effectués.
Le langage Rubic de base a été augmenté de fonctionnalités inspirées du
langage Ruby telles que la possibilité de commenter du code, \textbf{[à compléter au
fur et à mesure que l'on rajoute des fonctionnalités]}.

Ce projet a permis le développement et la mise en œuvre de nombreuses
compétences intéressantes telles que l'utilisation des outils Flex et Yacc, la
découverte du langage Ruby, l'initiation au langage LLVM et l'utilisation de
structures de données adaptées au typage des données, à la gestion des tables
de symboles et de la mémoire.

Sur les bases de ce qui a été accompli il serait possible de réaliser un
interpréteur pour le langage Rubic. Cela nécessiterai une modification des
structures employées afin de pouvoir, par exemple, associer des valeurs aux
variables. Les actions de génération de code effectuées lors du parsage
seraient alors remplacées par des calculs de valeur des expressions par
exemple. D'autres actions plus complexes devraient alors être développées.

\end{document}
